problem:  
Let \( M^n \subset \mathbb{R}^{n+1} \) be a complete, oriented hypersurface with induced metric \( g \) and second fundamental form \( A \). Assume that \( (M^n,g,f) \) satisfies the **almost gradient Ricci soliton** equation  
\[
\mathrm{Ric} + \nabla^2 f = \lambda(x) g,
\]
for smooth functions \( f, \lambda: M \to \mathbb{R} \). Define the **normalized traceless curvature** ratio  
\[
\Psi = \frac{\|A_0\|^2}{H^2},
\qquad A_0 = A - \tfrac{H}{n} g,
\]
where \(H\) is the mean curvature. Suppose \(M\) is complete, noncompact, has polynomial volume growth, and \( \Psi \) is constant on \(M\).  

**Question:**  
If \( \Psi \) is a strictly positive constant (i.e. \(M\) is not totally umbilic at any point), can such a hypersurface satisfy the almost Ricci soliton equation for a nonconstant soliton function \(f\)?  
Equivalently: do there exist non-umbilic, complete almost Ricci soliton hypersurfaces in Euclidean space with constant normalized traceless curvature ratio? If so, what structural constraints must the soliton potential \(f\) and \(\lambda(x)\) satisfy?  

Why is it a "good" problem:  
This problem strikes directly at the delicate interface between *extrinsic rigidity* and the *analytic structure* of almost Ricci solitons. The constancy of the traceless curvature ratio \( \Psi \) expresses that the hypersurface has globally uniform deviation from being umbilic—strong enough to suggest symmetry, yet flexible enough to avoid algebraic pinching reductions. Together with the differential almost soliton equation, this introduces a coupling between the potential function \(f\) and the geometry of how the shape operator’s spectrum distributes. Determining whether such constancy can persist under the analytic soliton constraints requires new interplay between elliptic PDE tools (Bochner identities, maximum principles in weighted manifolds) and classical submanifold geometry (Simons-type identities for \(A_0\)), thereby bridging *Ricci soliton analysis* and *extrinsic curvature classification*.  
The question is conceptually simple to pose but deeply unexplored; a solution would clarify whether almost soliton structures inherently prohibit uniformly non-umbilic geometries, thus shaping our understanding of rigidity within weighted extrinsic geometries.